\documentclass{beamer}
\usetheme{Madrid}
\usecolortheme{whale}
\usepackage[utf8]{inputenc}
\usepackage[spanish]{babel}
\usepackage{graphicx}
\usepackage{booktabs}

\title[Mega-Pipelines y CRISP-DM]{Día 08: Aplicaciones Industriales y Mega-Pipelines}
\subtitle{Diplomado en Ciencia de Datos}
\author{Docente Experto}
\institute{Universidad San Francisco Xavier (USFX)}
\date{Módulo 4: Aprendizaje Automático}

\begin{document}

\begin{frame}
    \titlepage
\end{frame}

\begin{frame}{Objetivos de la Sesión}
    \begin{itemize}
        \item \textbf{1.} Comprender la metodología \textbf{CRISP-DM} estándar en la industria.
        \item \textbf{2.} Identificar y prevenir el peligroso problema de la \textbf{Fuga de Datos (Data Leakage)}.
        \item \textbf{3.} Construir un \textbf{Mega-Pipeline} escalable utilizando Scikit-Learn.
        \item \textbf{4.} Realizar un análisis de impacto financiero a través de la \textbf{Matriz de Confusión}.
    \end{itemize}
\end{frame}

\begin{frame}{El Estándar Industrial: CRISP-DM}
    La metodología \textbf{Cross-Industry Standard Process for Data Mining} asegura que los proyectos de Inteligencia Artificial no fracasen en producción.
    \vspace{0.5cm}
    \begin{center}
        \includegraphics[height=5.5cm,keepaspectratio]{crisp_dm.png}
    \end{center}
\end{frame}

\begin{frame}{El Peligro Oculto: Data Leakage (Fuga de Datos)}
    \begin{alertblock}{Definición Fatal}
        Ocurre cuando el algoritmo accede a información del conjunto de pruebas (examen) durante la fase de entrenamiento.
    \end{alertblock}
    
    \textbf{¿Cómo sucede accidentalmente?}
    \begin{itemize}
        \item Imputar valores nulos (promedio/mediana) usando \textbf{todo} el dataset.
        \item Escalar variables utilizando el mínimo y máximo de \textbf{todo} el dataset.
    \end{itemize}
    \vspace{0.5cm}
    \textbf{Consecuencia:} El modelo reportará 99\% de exactitud en el laboratorio, pero fracasará catastróficamente el día que se lance a producción.
\end{frame}

\begin{frame}{La Solución Definitiva: El Pipeline}
    Para evitar la Fuga de Datos, empaquetamos las instrucciones de limpieza y el algoritmo en una sola tubería blindada.
    \vspace{0.5cm}
    \begin{center}
        \includegraphics[height=5.5cm,keepaspectratio]{pipeline_arch.png}
    \end{center}
\end{frame}

\begin{frame}{Evaluación Empresarial: Falsos Positivos vs Negativos}
    Las métricas de exactitud (\textit{Accuracy}) son inútiles para la Junta Directiva. Todo debe traducirse a \textbf{Dólares (Costos y Riesgos)}.
    
    \vspace{0.5cm}
    \begin{block}{Ejemplo Médico: Predicción de Cáncer}
        \begin{itemize}
            \item \textbf{Falso Positivo (Susto):} Decirle a un paciente sano que tiene cáncer. Costo: Ansiedad y exámenes adicionales.
            \item \textbf{Falso Negativo (Tragedia):} Decirle a un paciente con cáncer que está sano. Costo: Pérdida de la vida.
        \end{itemize}
    \end{block}
    
    \vspace{0.3cm}
    \textbf{Lección de Negocios:} Los Data Scientists ajustan los Pipelines para minimizar el tipo de error que le cueste más caro a la empresa.
\end{frame}

\begin{frame}{Siguientes Pasos}
    \begin{center}
        \Large \textbf{¡Manos al Código!}
        \vspace{0.5cm}
        
        \normalsize Pasemos al entorno de Google Colab para programar nuestro Mega-Pipeline desde cero y exportarlo a producción.
    \end{center}
\end{frame}

\end{document}
